\documentclass{article}
\usepackage[T1]{fontenc}
\usepackage{lmodern}
\usepackage{graphicx}
\usepackage{amsmath,amssymb}
\usepackage[a4paper,margin=2cm]{geometry}
\usepackage[absolute,overlay]{textpos}
\setlength{\TPHorizModule}{1mm}
\setlength{\TPVertModule}{1mm}
\usepackage[indonesian]{babel}
\usepackage{hyperref}
\setlength{\emergencystretch}{2em}
% unit: Halaman 44

\begin{document}

% === Halaman 44 (llm) ===

\section*{AES-256}

\begin{itemize}
    \item AES-256 dengan AES-128 hanya berbeda dalam hal:
    \begin{enumerate}
        \item Jumlah putaran 14 kali (AES-128 hanya 10 kali)
        \item Panjang kunci 256 bit (AES-128 hanya 128 bit)
        \item \textit{Expansion key} (untuk membangkitkan kunci putaran)
    \end{enumerate}

    \item \textit{Expansion key} pada AES-128 uritannya adalah sebagai berikut:
    
    RotWord $\rightarrow$ SubWord $\rightarrow$ XOR Rcon $\rightarrow$ XOR dengan word sebelumnya

    \item \textit{Expansion key} pada AES-256:
    
    Operasi SubWord tambahan pada word tertentu:
    
    jika $i \pmod 8 = 0$, dilakukan RotWord $\rightarrow$ SubWord $\rightarrow$ XOR Rcon
    
    jika $i \pmod 8 = 4$, dilakukan SubWord saja
\end{itemize}

\end{document}
